\subsection{Area between curves and a simple volume model}

If \(f(x)\ge g(x)\) on \([a,b]\), then the area between the graphs is
\[
\int_a^b \bigl(f(x)-g(x)\bigr)\,dx.
\]

\begin{examplebox}[Area between a line and a parabola]
On \([0,1]\), the line \(y=x\) lies above the parabola \(y=x^2\). Therefore
\[
A=\int_0^1 (x-x^2)\,dx
=\left[\frac{x^2}{2}-\frac{x^3}{3}\right]_0^1
=\frac16.
\]
\end{examplebox}

A simple volume of revolution about the horizontal axis is obtained by disks:
\[
V=\pi\int_a^b r(x)^2\,dx.
\]

\begin{examplebox}[Rotating a line segment]
Rotating \(y=x\) on \([0,1]\) about the horizontal axis produces a cone. The disk formula gives
\[
V=\pi\int_0^1 x^2\,dx=\frac{\pi}{3},
\]
which agrees with the elementary cone formula.
\end{examplebox}

\begin{interpretationbox}[Geometry as accumulation of slices]
Area accumulates vertical differences; volume accumulates cross-sectional areas. The integral changes meaning through the quantity being summed, not through a more complicated computational technique.
\end{interpretationbox}
